\documentclass[border=7bp]{standalone}

\usepackage{fontspec}
\usepackage{xeCJK}
\setmainfont{Calibri}
\setCJKmainfont{SimSun}[AutoFakeBold=1.5]
\usepackage{amsmath}
\usepackage{tikz}
\usetikzlibrary{arrows.meta}

\definecolor{Ink}{HTML}{25313D}
\definecolor{Blue}{HTML}{315E80}
\definecolor{PaleBlue}{HTML}{E7EFF5}
\definecolor{Muted}{HTML}{626B74}
\definecolor{PaleGray}{HTML}{F0F1F3}

\tikzset{
  >=Stealth,
  every node/.style={font=\fontsize{9.5bp}{12.6bp}\selectfont,text=Ink},
  cell/.style={draw=Ink,line width=.55pt,minimum size=7.2mm,inner sep=0pt},
  blank/.style={cell,fill=PaleGray},
  member/.style={cell,fill=PaleBlue,draw=Blue},
  tag/.style={font=\fontsize{9bp}{12bp}\selectfont,text=Muted},
  edge/.style={->,draw=Blue,line width=.7pt},
}

% 抽象状态：0 为空位，1..4 为组员，其余格画成浅点
\newcommand{\aboard}[3]{%
  \begin{scope}[shift={(#1,#2)}]
    \foreach \row [count=\r from 0] in {#3} {
      \foreach \v [count=\c from 0] in \row {
        \ifnum\v<0\relax
          \node[cell,text=black!25] at (\c*0.78,-\r*0.78) {$\cdot$};%
        \else
          \ifnum\v=0\relax
            \node[blank] at (\c*0.78,-\r*0.78) {};%
          \else
            \node[member] at (\c*0.78,-\r*0.78) {\v};%
          \fi
        \fi
      }
    }
  \end{scope}}

% Shared preamble for the report's editable vector figures.
\usepackage{fontspec}
\usepackage{xeCJK}
\setmainfont{Calibri}
\setCJKmainfont{SimSun}[AutoFakeBold=1.5]
\usepackage{amsmath}
\usepackage{tikz}
\usetikzlibrary{arrows.meta,positioning,calc,shapes.geometric,fit}
\definecolor{Ink}{HTML}{25313D}
\definecolor{Blue}{HTML}{315E80}
\definecolor{PaleBlue}{HTML}{E7EFF5}
\definecolor{Muted}{HTML}{626B74}
\definecolor{PaleGray}{HTML}{F0F1F3}
\tikzset{
  >=Stealth,
  every node/.style={font=\fontsize{10bp}{13bp}\selectfont,text=Ink},
  flow/.style={->,draw=Ink,line width=.65pt},
  frame/.style={rectangle,draw=Ink,line width=.65pt,align=center,
    minimum height=.85cm,text width=5cm,inner sep=5pt},
  emphasis/.style={frame,draw=Blue,fill=PaleBlue},
  supporting/.style={frame,draw=Muted,fill=PaleGray},
  annotation/.style={font=\fontsize{9bp}{12bp}\selectfont,text=Muted,align=center},
}

\begin{document}
\begin{tikzpicture}
  % 左侧：抽象状态棋盘
  \node[tag] at (1.17,0.62) {抽象状态（组员 $1,2,3,4$ 与空位）};
  \aboard{0}{0}{{1,-1,-1,-1},{-1,0,-1,2},{-1,-1,3,-1},{-1,4,-1,-1}}

  % 棋盘正下方标注位置序列，不与右侧内容争空间
  \node[tag,align=center] at (1.17,-3.95)
    {位置序列 $(5,\,0,\,7,\,10,\,13)$\\即（空位$,$ 1$,$ 2$,$ 3$,$ 4）};

  % 右侧：编码推导，纵向排在棋盘右侧的独立列，行距放宽
  \node[anchor=north west,align=left,text width=5.0cm] at (3.05,0.35)
  {\textbf{逐位变进制编码}\\[1.5bp]
   $r_0=5$\\
   $r_1=5\times15+0=75$\\
   $r_2=75\times14+5=1055$\\
   $r_3=1055\times13+7=13\,722$\\
   $r_4=13\,722\times12+9=164\,673$\\[4bp]
   序数 $=$ 该格在\textbf{未占用}格中的序号};

  % 数组示意整体放到下方，单独占一行
  \begin{scope}[shift={(0,-5.30)}]
    \draw[line width=.6pt] (0,0) rectangle (8.20,-0.85);
    \foreach \i in {1,...,7} {
      \draw[line width=.4pt] (\i*1.025,0) -- (\i*1.025,-0.85);
    }
    \draw[fill=PaleBlue,draw=Blue,line width=.65pt] (3.60,-0.02) rectangle (4.62,-0.83);
    \node[font=\fontsize{7.5bp}{9bp}\selectfont] at (4.11,-0.425) {$\mathbf{164\,673}$};
    \node[tag,anchor=west] at (8.40,-0.425) {\texttt{byte[]}，长 $524\,160$};
    \node[tag,anchor=east] at (-0.15,-0.425) {下标 0};
  \end{scope}

  \draw[edge] (4.11,-4.55) -- (4.11,-5.05);
  \node[tag,anchor=west] at (4.25,-4.80) {$\mathrm{rank}$ 即下标};

  \node[tag,align=center] at (4.10,-6.95)
    {每格一个字节，无符号读出即该组的距离分项；\\
     $16\cdot15\cdot14\cdot13\cdot12 = 524\,160$};
\end{tikzpicture}
\end{document}
